%% notes-tex/publishing-system/sections/5-bibliography.tex

\section{Bibliographies\secstatus{todo}}
\label{sec:bib}

\notesec{Two tiers}
The manuscript's \texttt{references.bib} comes from the group's \texttt{paper}
collection and nowhere else. That is the publication guarantee: every work the
paper cites lives in a library the whole group holds, so citations stay
recoverable after publication. A typeset note's \texttt{references.bib} is the
union of the group and personal collections of the same name as the document,
because technical notes cite more widely than the paper and much of that reading
lands in the personal library first; on a clean overlap the group entry wins.

\notesec{Two routes to the same file}
\texttt{make bib} exports through Better BibTeX on a machine with Zotero desktop
running, and \texttt{make bibcheck} reports the two collision cases first: one
citekey naming two different works, which is fatal because some citation would
point at the wrong paper, and one work under two citekeys, which only splits a
reference in the PDF. \texttt{make bib-pull} is the headless route: a nightly job
on the mirror host exports every declared bibliography over the Zotero Web API into
the mirror, pins each file's sha256 in a manifest, and tc-lit serves them. The pull
verifies the bytes against both the manifest and the response header and leaves
the target untouched on a mismatch. The two routes agree on citekeys and format
fields differently, so a switch between them shows as a diff.

\notesec{The gate on citations}
\texttt{make check} fails on any \texttt{\textbackslash cite} key that
\texttt{references.bib} does not contain. That is the mechanical form of the rule
that a paper is added to the library before it is cited, and it is the check that
would have caught a citation to a work nobody held.

\notesec{Who may write to the library}
Almost nobody, and no agent. The library is curated by hand. The one sanctioned
write is \texttt{zotero\_publish.py} uploading a document's own built PDF into that
document's output collection. Papers that are retrieved go to the mirror through
the literature subsystem, which reads Zotero and never writes to it. Adding a
reference is a curation decision that needs an explicit instruction each time;
noticing that one is missing is not that instruction. The reason is the loop
above: \texttt{make bib} regenerates the file from the library and \texttt{make
check} fails on any key the export no longer has, so both an unrequested addition
and its later cleanup change what a document is permitted to cite.
